\begin{frame}{Por que é um problema difícil?}
    Supondo que a verdadeira refletividade seja $r = [0,1,0,-0.5,0]$ 
    e a wavelet seja $w = [0.5,1,0.5]$;

    Dado que convolução produz $y = w * r$, em vez de vermos os dois picos de 
    refletividade diretamente, observamos versões deles "espalhadas" pela wavelet.

    Mas como saber se uma determinada estrutura em \(y\) veio da wavelet ou da 
    refletividade?

    Por exemplo, talvez uma wavelet mais estreita combinada com uma refletividade 
    mais larga produza praticamente o mesmo y.
\end{frame}